\chapter{El problema del diseño urbano óptimo}\label{cap:optimizacion}

\section{Conjunto de diseños}

Un diseño urbano es un par $d=(\mathcal S,\theta)$, donde $\mathcal S$ es un subconjunto de los cuatro
proyectos en decisión (extensión de la Línea 1 a la calle 100, L1X; Línea 2, L2; Línea 3, L3;
Regiotram del Norte, RTN) y $\theta\in\{0;\,0{,}5;\,1\}$ la intensidad de la política de proximidad y
policentrismo. El conjunto $\mathcal D$ tiene $2^4\times3=\NDisenos$ diseños. El diseño base,
$(\emptyset,0)$, contiene solo los proyectos comprometidos y sirve de referencia: el beneficio neto de
$d$ en el escenario $\omega$ es $BN(d,\omega)=\VP(\text{base},\omega)-\VP(d,\omega)$.

Los diseños se escriben en el documento y en el software como «\texttt{L3+RTN | θ=1}». Los años de
apertura de los proyectos en decisión son los de la \cref{tab:proyectos}; la \cref{cap:opciones}
levanta ese supuesto para la Línea 3.

\section{Formulación estocástica}

El problema de referencia es un programa estocástico con criterio media--CVaR,
\begin{equation}\label{eq:problema}
\min_{d\in\mathcal D}\;(1-w)\,\E\big[\VP(d,\cdot)\big]+w\,\CVaR_{0{,}95}\big[\VP(d,\cdot)\big]
\quad\text{s.\,a.}\quad\E\big[\mathrm{CAPEX}(d,\cdot)\big]\le B,
\end{equation}
donde $w\in[0,1]$ es el peso de la cola (aversión al riesgo) y $B$ un presupuesto de inversión
esperada. Con $w=0$ el planificador es neutral al riesgo; con $w=1$ minimiza el costo esperado en el
5\,\% de los peores escenarios.

\section{Aproximación por promedio muestral}

El valor esperado y el CVaR se reemplazan por sus versiones muestrales sobre los mismos $n$ escenarios
para todos los diseños (aproximación por promedio muestral, SAA). Como $\mathcal D$ es finito y pequeño,
el problema muestral se resuelve exactamente por enumeración: se evalúan los \NDisenos{} diseños, se
calculan $\hat\E$, $\widehat{\CVaR}$ y la probabilidad de beneficio neto positivo, y se identifica la
\emph{frontera de Pareto} en el plano $(\hat\E,\widehat{\CVaR})$: los diseños para los que no existe
otro con menor costo esperado y menor CVaR a la vez. Todo minimizador de \eqref{eq:problema} sin
restricción presupuestal está en esa frontera, y al recorrer $w$ de 0 a 1 se obtiene la ruta de diseños
óptimos.

\begin{remark}
La enumeración es posible porque cada evaluación es barata: el equilibrio de una red de 12 zonas se
vectoriza sobre los $n$ escenarios con productos matriciales $(n\times Z^2)\cdot(Z^2\times D)$, y los
\NDisenos{} diseños con 1.024 escenarios y 25 años se evalúan en pocos minutos en un computador
personal. Con muchos más proyectos, la enumeración debería sustituirse por una formulación entera mixta
con la representación lineal del CVaR de \citet{rockafellar2000}.
\end{remark}

\section{Errores de muestreo}

Cada estimador se reporta con su error estándar. Para el beneficio neto de un diseño frente a la base,
los números aleatorios comunes hacen que el error estándar sea mucho menor que el del costo de cada
diseño: con $n=\NEscenarios$, el error estándar del costo esperado del diseño base es \VPBaseEE{}
billones, mientras que el del beneficio neto del diseño óptimo es \BNOptimoEE.
